\section{Глава 7. n-шаговый бутстрэппинг}
\subsection{Упражнение 7.1(c. 179)}
В главе 6 мы отметили, что ошибку Монте-Карло можно записать в виде суммы TD-ошибок (6.6), если оценка ценности не изменяется от шага к шагу. Покажите, что $n$-шаговую ошибку в формуле (7.2) также можно записать в виде суммы TD-ошибок (опять-таки, если оценки ценности не изменяются), что обобщает ранее полученный результат.
\begin{equation}
    G_t - V(S_t) = \sum_{k=t}^{T-1} \gamma^{k-t} \delta_k.
    \tag{6.6}
\end{equation}
\begin{equation}
    V_{t+n}(S_t) \doteq V_{t+n-1}(S_t) + \alpha\left[G_{t:t+n} - V_{t+n-1}(S_t)\right], \quad 0 \le t < T.
    \tag{7.2}
\end{equation}
\begin{solution}
    $n$-шаговая ошибка в (7.2) --- это величина $G_{t:t+n} - V_{t+n-1}(S_t)$. По условию ценности не изменяются от шага к шагу, поэтому индекс времени у $V$ можно опустить: $V_{t+n-1}(S_t) = V(S_t)$. Нужно показать, что эта ошибка равна сумме $n$ подряд идущих TD-ошибок.

    Распишем $n$-шаговый доход через первое вознаграждение и оставшийся доход. Из определения (7.1) следует
    \[
    G_{t:t+n} = R_{t+1} + \gamma G_{t+1:t+n}.
    \]
    Это верно при $n \ge 1$. Здесь мы пользуемся тем, что усечённый доход с горизонтом $n$ шагов от момента $t$ --- это первое вознаграждение плюс усечённый доход с горизонтом $n-1$ шагов от момента $t+1$.

    Теперь действуем так же, как в выводе (6.6). Прибавим и вычтем $\gamma V(S_{t+1})$:
    \begin{align*}
        G_{t:t+n} - V(S_t) &= R_{t+1} + \gamma G_{t+1:t+n} - V(S_t) + \gamma V(S_{t+1}) - \gamma V(S_{t+1}) \\
        &= \bigl(R_{t+1} + \gamma V(S_{t+1}) - V(S_t)\bigr) + \gamma\bigl(G_{t+1:t+n} - V(S_{t+1})\bigr) \\
        &= \delta_t + \gamma\bigl(G_{t+1:t+n} - V(S_{t+1})\bigr).
    \end{align*}
    В скобке снова стоит $n$-шаговая ошибка, но уже для момента $t+1$ и с горизонтом на единицу меньше. Применим к ней то же самое разложение:
    \begin{align*}
        G_{t:t+n} - V(S_t) &= \delta_t + \gamma\bigl(G_{t+1:t+n} - V(S_{t+1})\bigr) \\
        &= \delta_t + \gamma \delta_{t+1} + \gamma^2 \bigl(G_{t+2:t+n} - V(S_{t+2})\bigr) \\
        &= \delta_t + \gamma \delta_{t+1} + \gamma^2 \delta_{t+2} + \dots + \gamma^{n-1}\bigl(G_{t+n:t+n} - V(S_{t+n})\bigr).
    \end{align*}
    Раскладывание останавливается, когда горизонт усечения становится нулевым. По определению $n$-шагового дохода $G_{t+n:t+n} = V(S_{t+n})$ (доход с нулевым числом шагов --- это просто оценка ценности). Поэтому последнее слагаемое обнуляется:
    \[
    \gamma^{n-1}\bigl(G_{t+n:t+n} - V(S_{t+n})\bigr) = \gamma^{n-1}\bigl(V(S_{t+n}) - V(S_{t+n})\bigr) = 0.
    \]
    В итоге получаем
    \[
    G_{t:t+n} - V(S_t) = \sum_{k=t}^{t+n-1} \gamma^{k-t} \delta_k.
    \]
    Это сумма $n$ подряд идущих TD-ошибок, что и требовалось.

    Замечание: при $n = T - t$ доход $G_{t:t+n}$ становится обыкновенным полным доходом $G_t$, верхний предел суммы превращается в $T-1$, и формула переходит в (6.6). Так что результат главы 6 --- это частный случай при $n \to T-t$.
\end{solution}

\subsection{Упражнение 7.2(c. 179)}
В случае $n$-шагового метода оценки ценности изменяются от шага к шагу,
поэтому алгоритм, в котором используется сумма TD-ошибок (см. предыдущее
упражнение) вместо ошибки в формуле (7.2), отличается от описанного выше.
Лучше он или хуже? Придумайте и запрограммируйте простой эксперимент,
чтобы ответить на этот вопрос эмпирически.
\begin{equation}
    V_{t+n}(S_t) \doteq V_{t+n-1}(S_t) + \alpha\left[G_{t:t+n} - V_{t+n-1}(S_t)\right],
    \quad 0 \le t < T.
    \tag{7.2}
\end{equation}
\begin{solution}
    В упражнении 7.1 равенство $G_{t:t+n}-V(S_t)=\sum_{k=t}^{t+n-1}\gamma^{k-t}\delta_k$
    выполнено точно только если все $\delta_k$ посчитаны с одной и той же $V$.
    В классическом алгоритме (7.2) бутстрап-член $V(S_{\tau+n})$ всегда свежий.
    В варианте с суммой TD-ошибок каждая $\delta_k$ считается один раз и
    сохраняется, поэтому разные ошибки используют разные версии $V$ ---
    обновление работает с устаревшими оценками.

    \textbf{Эксперимент.} Сеточный мир $5\times5$, завышенная инициализация
    $V(s)=+5$, метрика --- RMSE до $v_\pi$ по эпизодам, усреднённая по прогонам,
    $n\in\{1,2,4,8,16\}$.

    \begin{figure}[ht]
        \centering
        \includegraphics[width=\textwidth]{7_2_sol_rmse.png}
        \caption{RMSE классического алгоритма (7.2) и варианта с суммой TD-ошибок.}
    \end{figure}
    \begin{figure}[ht]
        \centering
        \includegraphics[width=0.7\textwidth]{7_2_sol.png}
        \caption{Разность RMSE: положительные значения --- классический точнее.}
    \end{figure}

    \textbf{Результаты.} При $n=1$ алгоритмы совпадают тождественно. В установившемся режиме разность положительна и растёт с $n$.
    В первые $\sim$30 эпизодов при больших $n$ разность уходит в минус: устаревшие
    ошибки тормозят резкий спад завышенной $V$, что временно ближе к истине.

    \textbf{Вывод.} Вариант с суммой сохранённых TD-ошибок \emph{хуже}: разница
    растёт с $n$ и исчезает при $n=1$. Причина --- использование устаревших оценок
    $V$; в коротком переходном режиме это может помогать, но в целом вредит.
\end{solution}
\notebook{07_02}{7.2}

\subsection{Упражнение 7.3(c. 181)}
Как вы думаете, почему в примерах из этой главы используется более масштабная задача о случайном блуждании (с 19, а не 5 состояниями)? Верно ли, что при меньшем блуждании более выгодным оказалось бы другое значение $n$? А что, если в большем блуждании изменить результат левой части с $0$ на $-1$? Повлияло бы это на выбор наилучшего значения?
\begin{solution}
    \textbf{Почему 19 состояний.} Цель рис. 7.2 --- показать, что лучше всего
    работают промежуточные $n$. Для этого нужно, чтобы крайние варианты были
    заметно хуже среднего. В короткой цепочке из 5 состояний эпизод слишком
    короткий: при $n$ больше длины эпизода $n$-шаговый доход превращается в
    полный доход, и все большие $n$ дают один и тот же метод Монте-Карло.
    Различимых значений $n$ остаётся мало, и выраженного минимума на кривой
    ошибки не получится. Длинная цепочка даёт достаточно места, чтобы развести
    методы по $n$ и увидеть U-образную кривую.
 
    \textbf{Другое $n$ для маленькой задачи.} Да. Оптимальное $n$ соразмерно
    длине эпизода. В цепочке из 5 состояний уже двух-трёх шагов хватает, чтобы
    обновление дошло до всех посещённых состояний, а большие $n$ схлопываются в
    Монте-Карло. Поэтому оптимум сместился бы к малым $n$ и стал бы менее
    выраженным.
 
    \textbf{Замена $0$ на $-1$ слева.} Да, повлияло бы. Выгода большого $n$ ---
    быстро доставлять коррекцию к удалённым состояниям; она велика там, где у
    этих состояний большая начальная ошибка. Замена награды меняет истинные
    ценности (диапазон $[0,1]$ переходит в $[-1,+1]$), а значит и профиль
    начальной ошибки $V_0 - v_\pi$ --- даже при одинаковой инициализации, так
    как $v_\pi$ у задач разные.
\end{solution}

\subsection{Упражнение 7.4(c. 183)}
Докажите, что $n$-шаговый доход в методе Sarsa (7.4) можно точно выразить
в терминах новой TD-ошибки:
\begin{multline}
    G_{t:t+n} = Q_{t-1}(S_t, A_t) \\
    + \sum_{k=t}^{\min(t+n,T)-1}
    \gamma^{k-t}\left[R_{k+1} + \gamma Q_k(S_{k+1}, A_{k+1}) - Q_{k-1}(S_k, A_k)\right].
    \tag{7.6}
\end{multline}
\begin{multline}
    G_{t:t+n} \doteq R_{t+1} + \gamma R_{t+2} + \dots + \gamma^{n-1} R_{t+n} \\
    + \gamma^n Q_{t+n-1}(S_{t+n}, A_{t+n}), \quad n \ge 1,\ 0 \le t < T-n.
    \tag{7.4}
\end{multline}
\begin{solution}
    Обозначим TD-ошибку для ценности действий
    \[
    \delta_k \doteq R_{k+1} + \gamma Q_k(S_{k+1}, A_{k+1}) - Q_{k-1}(S_k, A_k).
    \]
    Распишем сумму из правой части (7.6), считая для удобства $t+n \le T$
    (тогда верхний предел равен $t+n-1$):
    \begin{align*}
        \sum_{k=t}^{t+n-1} \gamma^{k-t} \delta_k
        &= \sum_{k=t}^{t+n-1} \gamma^{k-t} R_{k+1}
        + \sum_{k=t}^{t+n-1} \gamma^{k-t+1} Q_k(S_{k+1}, A_{k+1}) \\
        &\quad - \sum_{k=t}^{t+n-1} \gamma^{k-t} Q_{k-1}(S_k, A_k).
    \end{align*}
    В средней сумме сделаем замену индекса $j = k+1$, тогда она пробегает
    $j = t+1, \dots, t+n$:
    \[
    \sum_{k=t}^{t+n-1} \gamma^{k-t+1} Q_k(S_{k+1}, A_{k+1})
    = \sum_{j=t+1}^{t+n} \gamma^{j-t} Q_{j-1}(S_j, A_j).
    \]
    Это та же величина, что и последняя сумма, только с другими пределами.
    При вычитании совпадающие члены для $j = t+1, \dots, t+n-1$ сокращаются,
    остаются крайние:
    \[
    \begin{aligned}
    &\sum_{j=t+1}^{t+n} \gamma^{j-t} Q_{j-1}(S_j, A_j)
    - \sum_{k=t}^{t+n-1} \gamma^{k-t} Q_{k-1}(S_k, A_k) \\
    &\qquad = \gamma^n Q_{t+n-1}(S_{t+n}, A_{t+n}) - Q_{t-1}(S_t, A_t).
    \end{aligned}
    \]
    Подставляя обратно, получаем
    \[
    \sum_{k=t}^{t+n-1} \gamma^{k-t} \delta_k
    = \underbrace{\sum_{k=t}^{t+n-1} \gamma^{k-t} R_{k+1}
    + \gamma^n Q_{t+n-1}(S_{t+n}, A_{t+n})}_{=\ G_{t:t+n}\ \text{по (7.4)}}
    - Q_{t-1}(S_t, A_t).
    \]
    Первые два слагаемых --- это в точности определение $G_{t:t+n}$ из (7.4).
    Перенося $Q_{t-1}(S_t, A_t)$ влево, приходим к (7.6):
    \[
    G_{t:t+n} = Q_{t-1}(S_t, A_t) + \sum_{k=t}^{t+n-1} \gamma^{k-t} \delta_k.
    \]
    Если эпизод заканчивается раньше ($t+n > T$), доход обрывается на шаге $T$,
    и верхний предел суммы становится $\min(t+n,T)-1$ --- это и записано в (7.6).
\end{solution}

\subsection{Упражнение 7.5(c. 186)}
Напишите псевдокод описанного выше алгоритма предсказания ценности состояния
с разделённой стратегией.
\begin{solution}
    Алгоритм совпадает по структуре с $n$-шаговым TD-предсказанием, но
    $n$-шаговый доход вычисляется по рекуррентной формуле (7.13):
    \[
    G_{t:h} = \rho_t\bigl(R_{t+1} + \gamma G_{t+1:h}\bigr)
    + (1-\rho_t)V_{h-1}(S_t), \qquad \rho_t = \frac{\pi(A_t|S_t)}{b(A_t|S_t)},
    \]
    с базой $G_{h:h} = V_{h-1}(S_h)$. Рекурсия разворачивается в один цикл,
    идущий от $h$ к $\tau$: переменная $G$ инициализируется базой и на каждом
    шаге пересчитывается. Дополнительной памяти не требуется --- значения
    $\rho_k$, $R_{k+1}$, $S_k$ алгоритм уже хранит.

    \begin{algorithm}[H]
        \DontPrintSemicolon
        \KwIn{поведенческая стратегия $b$ такая, что $b(a|s) > 0$ для всех
        $s \in \mathcal{S},\ a \in \mathcal{A}$; целевая стратегия $\pi$}
        Параметры: размер шага $\alpha \in (0,1]$, целое $n > 0$\;
        Инициализировать $V(s)$ произвольно для всех $s \in \mathcal{S}$\;
        Все операции с $S_t$ и $R_t$ берут индекс по модулю $n+1$\;
        \BlankLine
        \ForEach{эпизод}{
            Инициализировать и сохранить $S_0 \ne$ заключительное состояние\;
            $T \leftarrow \infty$\;
            \For{$t = 0, 1, 2, \dots$}{
                \If{$t < T$}{
                    Выбрать действие $A_t \sim b(\cdot|S_t)$\;
                    Предпринять $A_t$, наблюдать и сохранить $R_{t+1}$, $S_{t+1}$\;
                    \If{$S_{t+1}$ --- заключительное состояние}{
                        $T \leftarrow t+1$\;
                    }
                }
                $\tau \leftarrow t - n + 1$ \tcp*{момент обновления}
                \If{$\tau \ge 0$}{
                    $h \leftarrow \min(\tau+n,\, T)$\;
                    $G \leftarrow V(S_h)$ \tcp*{база $G_{h:h}$}
                    \For{$k = h-1, h-2, \dots, \tau$}{
                        $\rho \leftarrow \dfrac{\pi(A_k|S_k)}{b(A_k|S_k)}$\;
                        $G \leftarrow \rho\,(R_{k+1} + \gamma G)
                        + (1-\rho)\,V(S_k)$\;
                    }
                    $V(S_\tau) \leftarrow V(S_\tau)
                    + \alpha\,[\,G - V(S_\tau)\,]$\;
                }
                \If{$\tau = T - 1$}{
                    \textbf{выйти из цикла}\;
                }
            }
        }
        \caption{$n$-шаговое предсказание $V \approx v_\pi$ с разделённой
        стратегией и переменным управлением}
    \end{algorithm}
\end{solution}

\subsection{Упражнение 7.6(c. 187)}
Докажите, что переменное управление в приведённых выше уравнениях не изменяет
математического ожидания дохода.
\begin{solution}
    Докажем для ценностей действий (формула 7.14); для ценностей состояний
    (7.13) рассуждение аналогично. Переменное управление --- это член,
    добавленный к доходу:
    \[
    C = \gamma\bar V_{h-1}(S_{t+1}) - \gamma\rho_{t+1}Q_{h-1}(S_{t+1}, A_{t+1}),
    \qquad \rho_{t+1} = \frac{\pi(A_{t+1}|S_{t+1})}{b(A_{t+1}|S_{t+1})}.
    \]
    Достаточно показать, что $\mathbb{E}_b[C] = 0$ --- тогда добавление $C$ не
    меняет математического ожидания дохода.

    Зафиксируем $S_{t+1} = s$ и возьмём условное математическое ожидание по
    выбору $A_{t+1} \sim b(\cdot|s)$. Первый член от $A_{t+1}$ не зависит и
    выходит из-под ожидания как есть. Второй член распишем суммой по действиям:
    \[
    \begin{aligned}
    \mathbb{E}_b\bigl[\rho_{t+1}Q_{h-1}(s, A_{t+1}) \mid s\bigr]
    &= \sum_a b(a|s)\,\frac{\pi(a|s)}{b(a|s)}\,Q_{h-1}(s, a) \\
    &= \sum_a \pi(a|s)\,Q_{h-1}(s, a).
    \end{aligned}
    \]
    Здесь $b(a|s)$ сократилось. По определению (7.8) ожидаемая ценность состояния
    есть усреднение $Q$ по целевой стратегии:
    \[
    \bar V_{h-1}(s) = \sum_a \pi(a|s)\,Q_{h-1}(s, a).
    \]
    Значит оба члена $C$ равны $\gamma\bar V_{h-1}(s)$, и
    \[
    \mathbb{E}_b[C \mid s]
    = \gamma\bar V_{h-1}(s) - \gamma\bar V_{h-1}(s) = 0
    \quad \text{при любом } s.
    \]
    По формуле полного математического ожидания $\mathbb{E}_b[C] = 0$.

    \textbf{Случай ценностей состояний (7.13).} Управляющий член ---
    $V_{h-1}(S_t) - \rho_t V_{h-1}(S_t)$. Оценка $V_{h-1}(S_t)$ от действия не
    зависит, поэтому, используя известный факт $\mathbb{E}_b[\rho_t \mid s] = 1$,
    \[
    \mathbb{E}_b\bigl[V_{h-1}(s) - \rho_t V_{h-1}(s) \mid s\bigr]
    = V_{h-1}(s)\bigl(1 - \mathbb{E}_b[\rho_t \mid s]\bigr)
    = V_{h-1}(s)(1 - 1) = 0.
    \]

    В обоих случаях работает одно и то же сокращение $b(a|s)\,\rho = \pi(a|s)$:
    для состояний оно сворачивается до $\mathbb{E}_b[\rho] = 1$, для действий ---
    до $\mathbb{E}_b[\rho Q] = \bar V$. Управляющий член при этом не константа:
    он коррелирует с доходом и снижает дисперсию, не внося смещения.
\end{solution}

\subsection{*Упражнение 7.7(c. 192)}
Напишите псевдокод описанного выше алгоритма предсказания ценности действий
с разделённой стратегией. Особое внимание обратите на условия завершения
рекурсии при достижении горизонта или конца эпизода.
\begin{solution}
    Алгоритм совпадает по структуре с $n$-шаговым предсказанием, но
    $n$-шаговый доход вычисляется по рекуррентной формуле (7.14):
    \[
    G_{t:h} = R_{t+1} + \gamma\rho_{t+1}\bigl(G_{t+1:h} - Q_{h-1}(S_{t+1}, A_{t+1})\bigr)
    + \gamma\bar V_{h-1}(S_{t+1}),
    \]
    где $\rho_{t+1} = \pi(A_{t+1}|S_{t+1})/b(A_{t+1}|S_{t+1})$ и
    $\bar V(s) = \sum_a \pi(a|s)Q(s,a)$. Рекурсия разворачивается в один цикл
    от $h$ к $\tau$. Условие завершения зависит от того, оборвался эпизод или
    нет:
    \begin{itemize}
        \item если $h \le T$ (горизонт внутри эпизода), база --- $G_{h:h} = Q(S_h, A_h)$;
        \item если $h > T$ (эпизод закончился раньше), база --- $G_{T-1:h} = R_T$.
    \end{itemize}

    \begin{algorithm}[H]
        \small
        \DontPrintSemicolon
        \KwIn{поведенческая стратегия $b$ с $b(a|s) > 0$ для всех $s, a$;
        целевая стратегия $\pi$}
        Параметры: размер шага $\alpha \in (0,1]$, целое $n > 0$\;
        Инициализировать $Q(s,a)$ произвольно для всех $s \in \mathcal{S},\ a \in \mathcal{A}$\;
        Все операции с $S_t$, $A_t$, $R_t$ берут индекс по модулю $n+1$\;
        \BlankLine
        \ForEach{эпизод}{
            Инициализировать $S_0 \ne$ заключительное состояние\;
            Выбрать действие $A_0 \sim b(\cdot|S_0)$\;
            $T \leftarrow \infty$\;
            \For{$t = 0, 1, 2, \dots$}{
                \If{$t < T$}{
                    Предпринять $A_t$, наблюдать и сохранить $R_{t+1}$, $S_{t+1}$\;
                    \eIf{$S_{t+1}$ --- заключительное состояние}{
                        $T \leftarrow t+1$\;
                    }{
                        Выбрать и сохранить $A_{t+1} \sim b(\cdot|S_{t+1})$\;
                    }
                }
                $\tau \leftarrow t - n + 1$ \tcp*{момент обновления}
                \If{$\tau \ge 0$}{
                    $h \leftarrow \tau + n$ \tcp*{горизонт}
                    \BlankLine
                    \tcc{База рекурсии — два случая}
                    \eIf{$h \ge T$}{
                        $G \leftarrow R_T$ \tcp*{эпизод закончился: $G_{T-1:h}=R_T$}
                        $k_{\text{start}} \leftarrow T - 2$\;
                    }{
                        $G \leftarrow Q(S_h, A_h)$ \tcp*{горизонт внутри: $G_{h:h}=Q(S_h,A_h)$}
                        $k_{\text{start}} \leftarrow h - 1$\;
                    }
                    \BlankLine
                    \tcc{Разворачиваем рекурсию (7.14) от $k_{\text{start}}$ к $\tau$}
                    \For{$k = k_{\text{start}},\ k_{\text{start}}-1,\ \dots,\ \tau$}{
                        $\rho \leftarrow \dfrac{\pi(A_{k+1}|S_{k+1})}{b(A_{k+1}|S_{k+1})}$\;
                        $\bar V \leftarrow \sum_a \pi(a|S_{k+1})\,Q(S_{k+1}, a)$\;
                        $G \leftarrow R_{k+1}
                        + \gamma\rho\,\bigl(G - Q(S_{k+1}, A_{k+1})\bigr)
                        + \gamma\bar V$\;
                    }
                    $Q(S_\tau, A_\tau) \leftarrow Q(S_\tau, A_\tau)
                    + \alpha\,[\,G - Q(S_\tau, A_\tau)\,]$\;
                }
            }
            \If{$\tau = T - 1$}{
                \textbf{выйти из цикла}\;
            }
        }
        \caption{$n$-шаговое предсказание $Q \approx q_\pi$ с разделённой
        стратегией и переменным управлением}
    \end{algorithm}

    \textbf{Пояснение к условиям завершения.} Если горизонт $h$ не достигнут
    раньше конца эпизода ($h < T$), доход обрывается на оценке $Q(S_h, A_h)$ и
    цикл идёт от $h-1$. Если же эпизод закончился до горизонта ($h \ge T$),
    последнего действия $A_h$ просто не существует, поэтому базой служит
    последнее реальное вознаграждение $R_T$, а цикл начинается с $T-2$. В обоих
    случаях рекурсия сворачивается в один проход без отдельной рекурсивной
    функции.
\end{solution}

\subsection{Упражнение 7.8(c. 187)}
Покажите, что общий (с разделённой стратегией) вариант $n$-шагового дохода
(7.13) всё ещё можно точно и компактно записать в виде суммы основанных
на состояниях TD-ошибок (6.5), если приближённая функция ценности состояний
не изменяется.
\begin{equation}
    G_{t:h} \doteq \rho_t\left(R_{t+1} + \gamma G_{t+1:h}\right) + (1-\rho_t)V(S_t),
    \quad t < h < T.
    \tag{7.13}
\end{equation}
\begin{equation}
    \delta_t \doteq R_{t+1} + \gamma V(S_{t+1}) - V(S_t).
    \tag{6.5}
\end{equation}
\begin{solution}
    Раз $V$ не меняется, пишем её без индексов. Вынесем $V(S_t)$ из (7.13) и
    вставим в скобку $\gamma V(S_{t+1})$, чтобы выделить $\delta_t$:
    \begin{align*}
        G_{t:h} &= V(S_t) + \rho_t\left(R_{t+1} + \gamma G_{t+1:h} - V(S_t)\right)\\
                &= V(S_t) + \rho_t\left(\delta_t + \gamma\left(G_{t+1:h} - V(S_{t+1})\right)\right).
    \end{align*}

    Обозначим $E_t \doteq G_{t:h} - V(S_t)$. Тогда $E_t = \rho_t(\delta_t + \gamma E_{t+1})$,
    причём из условия завершения $G_{h:h} = V(S_h)$ следует $E_h = 0$. Раскручивая
    рекурсию вниз:
    \begin{equation*}
        E_t = \sum_{k=t}^{h-1} \gamma^{k-t}\left(\prod_{i=t}^{k}\rho_i\right)\delta_k,
    \end{equation*}
    откуда
    \begin{equation}
        G_{t:h} = V(S_t) + \sum_{k=t}^{h-1}
        \gamma^{k-t}\left(\prod_{i=t}^{k}\rho_i\right)\delta_k.
        \tag{$\ast$}
    \end{equation}

    Проверим края: при всех $\rho_i = 1$ получаем формулу (6.6), а при $\rho_t = 0$
    --- $G_{t:h} = V(S_t)$, что совпадает с (7.13). Равенство ($\ast$) точное только
    потому, что $V$ неизменна: иначе члены $\gamma V(S_{t+1})$ на соседних шагах
    брались бы из разных версий $V$ и не сокращались бы.
\end{solution}

\subsection{Упражнение 7.9(c. 187)}
Повторите предыдущее упражнение для варианта $n$-шагового дохода с разделённой
стратегией для ценностей действий (7.14) и TD-ошибки в смысле Expected Sarsa
(величина в квадратных скобках в формуле (6.9)).
\begin{equation}
    G_{t:h} \doteq R_{t+1} + \gamma\left(\rho_{t+1}\left(G_{t+1:h} - Q(S_{t+1}, A_{t+1})\right)
    + \bar{V}(S_{t+1})\right), \quad t < h \le T.
    \tag{7.14}
\end{equation}
\begin{equation}
    \delta_t \doteq R_{t+1} + \gamma\bar{V}(S_{t+1}) - Q(S_t, A_t),
    \qquad \bar{V}(s) \doteq \sum_a \pi(a\,|\,s)\,Q(s, a).
    \tag{6.9}
\end{equation}
\begin{solution}
    Раз $Q$ не меняется, пишем её без индексов. В (7.14) первое действие не
    взвешивается по значимости, поэтому рекурсия здесь по $Q$, а не по $V$.
    Вынесем $Q(S_t, A_t)$ и подставим $\gamma\bar V(S_{t+1})$, чтобы выделить $\delta_t$:
    \begin{align*}
        G_{t:h}
        &= Q(S_t, A_t)
        + \underbrace{R_{t+1} + \gamma\bar V(S_{t+1}) - Q(S_t, A_t)}_{\delta_t}\\
        &\quad + \gamma\rho_{t+1}\left(G_{t+1:h} - Q(S_{t+1}, A_{t+1})\right)\\
        &= Q(S_t, A_t) + \delta_t
        + \gamma\rho_{t+1}\left(G_{t+1:h} - Q(S_{t+1}, A_{t+1})\right).
    \end{align*}

    Обозначим $E_t \doteq G_{t:h} - Q(S_t, A_t)$. Тогда $E_t = \delta_t + \gamma\rho_{t+1}E_{t+1}$,
    причём из условия завершения рекурсии $E_h = 0$. Обрати внимание: первый член
    $\delta_t$ идёт без коэффициента $\rho$, а дальше каждый следующий уровень
    домножается на $\rho_{t+1}, \rho_{t+2}, \dots$. Раскручивая:
    \begin{equation*}
        E_t = \sum_{k=t}^{h-1} \gamma^{k-t}\left(\prod_{i=t+1}^{k}\rho_i\right)\delta_k,
    \end{equation*}
    где при $k=t$ пустое произведение равно $1$. Отсюда
    \begin{equation}
        G_{t:h} = Q(S_t, A_t) + \sum_{k=t}^{h-1}
        \gamma^{k-t}\left(\prod_{i=t+1}^{k}\rho_i\right)\delta_k.
        \tag{$\ast$}
    \end{equation}

    Отличие от упражнения 7.8 в том, что произведение коэффициентов начинается с
    $i=t+1$: первое действие $A_t$ уже выбрано и в выборку по значимости не входит.
    Как и там, равенство ($\ast$) точное лишь потому, что $Q$ неизменна --- иначе
    члены $\gamma\bar V(S_{t+1})$ на соседних шагах брались бы из разных версий $Q$
    и не сокращались бы.
\end{solution}

\subsection{Упражнение 7.10(c. 187)}
Придумайте не очень сложную задачу предсказания
с разделённой стратегией и воспользуйтесь ей, чтобы показать, что алгоритм
обучения с разделённой стратегией, в котором используются уравнения (7.13) и
(7.2), эффективнее более простого алгоритма, в котором используются уравнения
(7.1) и (7.9).
\begin{equation}
    \begin{aligned}
    G_{t:t+n} \doteq{}& R_{t+1} + \gamma R_{t+2} + \cdots
    + \gamma^{n-1}R_{t+n} \\
    &+ \gamma^n V_{t+n-1}(S_{t+n}), \quad n \ge 1,\ 0 \le t < T-n.
    \end{aligned}
    \tag{7.1}
\end{equation}
\begin{equation}
    V_{t+n}(S_t) \doteq V_{t+n-1}(S_t) + \alpha\left[G_{t:t+n} - V_{t+n-1}(S_t)\right],
    \quad 0 \le t < T.
    \tag{7.2}
\end{equation}
\begin{equation}
    V_{t+n}(S_t) \doteq V_{t+n-1}(S_t)
    + \alpha\rho_{t:t+n-1}\left[G_{t:t+n} - V_{t+n-1}(S_t)\right],
    \quad 0 \le t < T.
    \tag{7.9}
\end{equation}
\begin{equation}
    G_{t:h} \doteq \rho_t\left(R_{t+1} + \gamma G_{t+1:h}\right) + (1-\rho_t)V_{h-1}(S_t),
    \quad t < h < T.
    \tag{7.13}
\end{equation}
\begin{solution}
    \textbf{Код смотри в отдельном файле репозитория.}
    \textbf{Задача.} Случайное блуждание на 19 состояниях (пример 7.1), награда
    $+1$ справа и $-1$ слева. Поведенческая стратегия $b$ равновероятная;
    целевая $\pi$ задаётся одним числом $\pi(\text{вправо}\,|\,s)=p$ во всех
    состояниях, так что $\rho$ на шаге равен $p/0{,}5$ для хода вправо и
    $(1-p)/0{,}5$ для хода влево. Берём $n=4$ и $p\in\{0{,}1;\,0{,}3;\,0{,}5\}$.
    Истинные $v_\pi$ для каждого $p$ считаются точно из уравнений Беллмана,
    метрика --- RMSE до $v_\pi$ по состояниям и эпизодам, усреднённый по 100
    прогонам. При $p=0{,}5$ задача on-policy, поэтому оба метода обязаны совпасть
    (это проверка кода; при $n=1$ они совпадают тождественно).

    \begin{figure}[ht]
        \centering
        \includegraphics[width=\textwidth]{alpha_0_1.png}
        \caption{$n=4$, $\alpha=0{,}1$: при малом шаге методы почти неотличимы.}
    \end{figure}
    \begin{figure}[ht]
        \centering
        \includegraphics[width=\textwidth]{alpha_0_3.png}
        \caption{$n=4$, $\alpha=0{,}3$: при $p=0{,}1$ обычный метод (7.9)
        разваливается, control variate остаётся устойчивым (выглядит как ноль
        из-за масштаба оси).}
    \end{figure}
    \begin{figure}[H]
        \centering
        \includegraphics[width=0.7\textwidth]{log_alpha_0_3.png}
        \caption{Тот же случай $p=0{,}1$ в лог-масштабе: обе кривые видны
        раздельно, control variate точнее обычного на два порядка.}
    \end{figure}

    \textbf{Результаты.} При $p=0{,}5$ кривые совпадают --- веса равны 1. С ростом
    смещения $\pi$ коэффициенты $\rho$ всё сильнее отличаются от 1, и на окне
    $n=4$ их произведения разлетаются в широком диапазоне. При малом шаге
    $\alpha=0{,}1$ это почти незаметно. При $\alpha=0{,}3$ и сильном off-policy
    $p=0{,}1$ обычный метод (7.9) умножает обновление на огромные произведения
    $\rho$ и расходится (RMSE доходит до сотен), тогда как control variate
    усекает доход до текущей оценки $V$ в точках с малым $\rho$ и потому остаётся
    устойчивым.

    \textbf{Вывод.} Алгоритм с control variate (7.13)+(7.2) \emph{эффективнее}:
    при сильном расхождении $\pi$ и $b$ он точнее и не разваливается при больших
    шагах, а при $n=1$ и в on-policy случае совпадает с обычным.
\end{solution}
\notebook{07_10}{7.10}

\subsection{Упражнение 7.11(c. 187)}
Покажите, что если приближённые ценности действий не изменяются, то доход в
алгоритме обновления по дереву (7.16) можно записать в виде суммы основанных
на математическом ожидании TD-ошибок:
\begin{equation*}
    G_{t:t+n} = Q(S_t, A_t)
    + \sum_{k=t}^{\min(t+n,\,T-1)} \delta_k \prod_{i=t+1}^{k} \gamma\pi(A_i\,|\,S_i),
\end{equation*}
где $\delta_t \doteq R_{t+1} + \gamma\bar{V}_t(S_{t+1}) - Q(S_t, A_t)$, а $\bar V_t$
определено формулой (7.8).
\begin{equation}
    \begin{aligned}
    G_{t:t+n} \doteq{}& R_{t+1}
    + \gamma\sum_{a\ne A_{t+1}}\pi(a\,|\,S_{t+1})Q(S_{t+1},a)\\
    &+ \gamma\pi(A_{t+1}\,|\,S_{t+1})G_{t+1:t+n}, \quad t < T-1.
    \end{aligned}
    \tag{7.16}
\end{equation}
\begin{solution}
    Раз $Q$ не меняется, пишем её без индексов. Сумма по $a\ne A_{t+1}$ -- это
    полное $\bar V(S_{t+1})$ без члена выбранного действия:
    \begin{equation*}
        \sum_{a\ne A_{t+1}}\pi(a|S_{t+1})Q(S_{t+1},a)
        = \bar V(S_{t+1}) - \pi(A_{t+1}|S_{t+1})Q(S_{t+1},A_{t+1}).
    \end{equation*}
    Подставим это в (7.16) и сгруппируем слагаемые с $\gamma\pi(A_{t+1}|S_{t+1})$:
    \begin{align*}
        G_{t:t+n} &= \underbrace{R_{t+1} + \gamma\bar V(S_{t+1})}_{\delta_t + Q(S_t,A_t)}
                     + \gamma\pi(A_{t+1}|S_{t+1})\bigl(G_{t+1:t+n} - Q(S_{t+1},A_{t+1})\bigr)\\
                  &= Q(S_t,A_t) + \delta_t
                     + \gamma\pi(A_{t+1}|S_{t+1})\bigl(G_{t+1:t+n} - Q(S_{t+1},A_{t+1})\bigr).
    \end{align*}

    Обозначим $E_t \doteq G_{t:t+n} - Q(S_t, A_t)$. Тогда
    \begin{equation*}
        E_t = \delta_t + \gamma\pi(A_{t+1}|S_{t+1})\,E_{t+1}.
    \end{equation*}
    В отличие от упражнений 7.8--7.9 множитель при рекурсивном члене здесь не
    $\gamma\rho$, а $\gamma\pi(A_{t+1}|S_{t+1})$ -- это и даёт произведение
    $\prod\gamma\pi(A_i|S_i)$. Раскручивая рекурсию:
    \begin{equation*}
        E_t = \sum_{k=t}^{\min(t+n,\,T-1)} \delta_k \prod_{i=t+1}^{k}\gamma\pi(A_i|S_i),
    \end{equation*}
    где при $k=t$ произведение пустое и равно $1$.

    Завершение рекурсии. Если окно достаёт до конца эпизода ($t+n\ge T$), то
    последний шаг $\tau=T-1$ даёт $G_{T-1:t+n} = R_T$, и поскольку $S_T$
    терминально, $\bar V(S_T)=0$, так что $E_{T-1}=\delta_{T-1}$ точно --- отсюда
    верхний предел $\min(t+n,T-1)$. Если же окно закрывается раньше ($t+n<T$),
    отдельного хвостового члена $E_{t+n}$ не возникает: в последней точке доход
    становится одношаговым (7.15), для которого $G-Q=\delta$ по определению.
    Проверка при $n=1$: $G_{t:t+1}=R_{t+1}+\gamma\bar V(S_{t+1})$, то есть
    $E_t=\delta_t$ --- ровно один член суммы с пустым произведением.

    Возвращаясь к $G_{t:t+n} = Q(S_t,A_t) + E_t$, получаем требуемое равенство.
\end{solution}
